\documentclass[11pt,a4paper]{article}
\usepackage[utf8]{inputenc}
\usepackage[T1]{fontenc}
\usepackage{amsmath,amssymb}
\usepackage{booktabs}
\usepackage{hyperref}
\usepackage[margin=1in]{geometry}
\usepackage{natbib}

\title{\textbf{Ethics Packs:}\\Knowledge Graph Injection for Moral Reasoning in Language Models\\{\large A Pilot Study with a Measurement Confound}}
\author{Nexus\textsuperscript{1} \and Lyra\textsuperscript{1} \and Thomas Edrington\textsuperscript{1}\\
\textsuperscript{1}Liberation Labs / Transparent Humboldt Coalition}
\date{May 2026}

\begin{document}
\maketitle

\begin{abstract}
Theory-matched ethics knowledge graphs improve moral reasoning scores in language models: Aristotelian packs on Aristotelian dilemmas produce the largest gain (+0.074, from 0.572 to 0.646), with four of five philosophical frameworks showing positive deltas. We construct 30 knowledge graphs from the Stanford Encyclopedia of Philosophy, encode their structure as walk topology (random-walk transition matrices that capture which concepts connect to which), and inject them as text context into a 30B Mixture-of-Experts model. The overall effect across frameworks is +0.038 ($n{=}5$ per framework, $p{=}0.193$, Cohen's $d{=}0.268$). Because the evaluation uses keyword matching, the measured improvement could reflect vocabulary priming rather than genuine reasoning; an LLM-judge re-evaluation is needed to distinguish these explanations. The methodology---automated ethics graph construction, walk encoding, and theory-matched injection---is validated; the measurement requires follow-up.
\end{abstract}

\section{Introduction}

Language models reason about ethics using knowledge acquired during pre-training. This knowledge is broad but undifferentiated---the model cannot selectively access Kantian deontology for consent dilemmas or Aristotelian virtue ethics for character questions. The knowledge is baked in, not routed.

Knowledge Packs \citep{pustovit2026knowledge} demonstrated that domain knowledge can be pre-computed as KV cache state and injected at inference time, adding expertise without retraining. Our prior work \citep{nexus2026topology} extended this from text content to graph topology, showing that walk-encoded knowledge graphs---where random-walk transition matrices capture the connectivity structure of a knowledge domain---enable language models to recover structural relationships. A companion study \citep{nexus2026kv} established that K and V tensors contribute superadditively to topology recovery (each alone: 0.333; combined: 1.000).

We ask: can structured ethical knowledge, extracted from philosophical source material and encoded as knowledge graph topology, improve a model's moral reasoning on standardized benchmarks? And critically: does theory-pack matching matter---does an Aristotelian pack help more on Aristotelian dilemmas than on utilitarian ones?

\section{Methods}

\subsection{Ethics Knowledge Graphs}

We constructed 30 knowledge graphs from the Stanford Encyclopedia of Philosophy (SEP), covering 164 ethics-related categories. For each category, we processed 20 entries through an Open Information Extraction (OpenIE) pipeline---an automated system that identifies entity-relationship triples (e.g., ``autonomy \textrightarrow{} requires \textrightarrow{} consent'') from unstructured text---using DeepSeek v2 16B as the extraction model. From the extracted triples, we built a NetworkX graph and computed walk encodings (random walk transition matrix accumulated over 5 steps).

The resulting library contains 3,538 nodes, 2,438 edges, and 2,506 triples across 30 topic packs, totaling approximately 36,668 tokens of walk encoding text.

\subsection{Theory-Pack Matching}

We mapped five MoReBench ethical frameworks to SEP packs:

\begin{table}[h]
\centering
\begin{tabular}{ll}
\toprule
\textbf{Framework} & \textbf{Matched Packs} \\
\midrule
Aristotelian Virtue Ethics & aristotle-ethics, ethics-ancient, moral-character \\
Kantian Deontology & autonomy-moral, informed-consent, personal-autonomy \\
Act Utilitarianism & moral-cognitivism, reasoning-moral \\
Scanlonian Contractualism & justice-climate, civil-rights \\
Gauthierian Contractarianism & ethics-ai, computing-responsibility \\
\bottomrule
\end{tabular}
\caption{Theory-to-pack mapping. Packs selected by conceptual alignment with each framework's core concerns.}
\label{tab:mapping}
\end{table}

\subsection{Evaluation}

\textbf{Benchmark.} MoReBench \citep{morebench2025}, a standardized moral reasoning benchmark that presents ethical dilemmas and scores responses against rubric criteria for each philosophical framework: 150 theory-specific dilemmas (30 per framework) with 15,285 weighted rubric criteria across five evaluation dimensions: identifying moral issues, logical reasoning, clear articulation, helpful outcomes, and harmless outcomes.

\textbf{Design.} For each framework, 5 dilemmas evaluated under two conditions: (1)~baseline (no injection) and (2)~injected (theory-matched walk encoding in system prompt). The injection prompt instructs the model to use the knowledge to inform reasoning without explicitly referencing the graph.

\textbf{Scoring.} Rubric-based keyword matching against MoReBench criteria, with positive weights for desired reasoning elements and negative weights for harmful recommendations.

\textbf{Model.} Qwen3-30B-A3B (Mixture of Experts, 30.5B total, $\sim$3B active per token) via Ollama, temperature 0.3, 4000 max tokens, 600-second timeout.

\subsection{Timeout Confound}

An initial run (v3) used a 300-second timeout. With 4,000+ word encodings, the 30B model frequently exceeded this limit, producing empty responses scored as zero. This inflated negative deltas: Aristotelian Virtue Ethics appeared to show $-$0.236 when the true effect (v4, 600s timeout, zero timeouts) is $+$0.074. We report both runs for methodological transparency.

\section{Results}

\begin{table}[h]
\centering
\begin{tabular}{lccr}
\toprule
\textbf{Framework} & \textbf{Baseline} & \textbf{Injected} & \textbf{$\Delta$} \\
\midrule
Aristotelian Virtue Ethics & 0.572 & \textbf{0.646} & \textbf{+0.074} \\
Gauthierian Contractarianism & 0.403 & 0.470 & +0.067 \\
Kantian Deontology & 0.480 & 0.537 & +0.058 \\
Act Utilitarianism & 0.487 & 0.513 & +0.026 \\
Scanlonian Contractualism & 0.563 & 0.530 & $-$0.033 \\
\midrule
\textbf{Overall} & \textbf{0.501} & \textbf{0.539} & \textbf{+0.038} \\
\bottomrule
\end{tabular}
\caption{MoReBench rubric scores by framework. Four of five frameworks improve with theory-matched pack injection. Zero timeouts in this run.}
\label{tab:results}
\end{table}

The best-matched condition (Aristotelian packs on Aristotelian dilemmas) produces the strongest improvement (+0.074). The overall effect across all frameworks is +0.038. Scanlonian Contractualism shows a slight negative ($-$0.033), within noise for 5 dilemmas.

\section{Discussion}

\subsection{What This Establishes}

The methodology works: automated construction of ethics knowledge graphs from SEP, walk encoding, and context injection at inference time produce a functioning pipeline. Theory-pack matching produces the largest deltas (+0.074 for Aristotle-on-Aristotle), and four of five frameworks show positive effects. The pipeline enables targeted ethical knowledge routing without retraining.

\subsection{What Remains Preliminary}

The overall +0.038 effect is not statistically significant ($n{=}25$, paired $t$-test $p{=}0.193$, Cohen's $d{=}0.268$). More critically, the keyword-based evaluation cannot distinguish genuine reasoning improvement from vocabulary priming: injecting ethics-specific terminology into the context causes the model to adopt that vocabulary, which the keyword scorer counts as hits. A vocabulary-control condition (same token length, shuffled text, no graph structure) or an LLM-as-judge evaluation is required to determine whether the improvement reflects structural reasoning scaffolding or surface-level vocabulary adoption.

\subsection{Theory Matching Matters}

The strongest improvement occurs when the injected pack matches the dilemma's framework (+0.074 for Aristotle-on-Aristotle). This suggests that the knowledge graph provides framework-specific conceptual relationships that the model leverages for reasoning within that framework. For production systems, this implies a routing architecture: identify the relevant ethical framework, select the matched pack, inject at inference time.

\subsection{Walk Encoding as Ethical Scaffolding}

Walk encoding captures topology---which concepts connect to which, with what strength---but not semantic content (the arguments, the reasoning, the philosophical positions). Despite this, it improves reasoning quality. We hypothesize that the topology provides \textit{conceptual scaffolding}: the model already knows what ``autonomy'' and ``consent'' mean from pre-training, but the injected topology tells it that these concepts are structurally related to ``dignity'' and ``self-determination'' in this specific ethical framework. The graph doesn't teach content; it teaches structure.

\subsection{Limitations}

\textbf{Sample size.} Five dilemmas per framework. Larger samples are needed for statistical significance claims.

\textbf{Keyword scoring.} MoReBench rubric scoring uses keyword matching, which undercounts semantically correct but differently-worded responses. LLM-judge scoring would provide more accurate evaluation.

\textbf{Single model.} Results from one model (Qwen3-30B-A3B MoE). Cross-model replication needed.

\textbf{Walk encoding only.} Graph topology without semantic content. Richer encodings (triples text, source excerpts) might produce stronger effects.

\textbf{Pack matching heuristic.} Theory-to-pack mapping was manual. Automated routing (e.g., embedding similarity) would test whether matching can be learned.

\section{Future Work}

\textbf{Pack routing.} Automated selection of the most relevant ethics pack based on query embedding similarity. The routing hypothesis: right pack helps, wrong pack hurts less than no pack, optimal routing maximizes benefit.

\textbf{Semantic enrichment.} Inject triple text or source excerpts alongside walk encoding. Topology provides structure; text provides arguments. Both together may be superadditive, as K and V tensors are for cache injection \citep{nexus2026kv}.

\textbf{Circular emotion geometry.} Recent work \citep{sun2026circumplex} establishes that emotion in LLMs occupies a 2D circular subspace (Russell's circumplex). Ethics packs that encode emotional valence and arousal alongside conceptual topology could modulate both reasoning structure and affective engagement.

\textbf{Oracle integration.} Deploy ethics packs as a runtime module in the Oracle ethical reasoning system, with SIRA-based routing selecting the appropriate pack per query.

\section{Conclusion}

Ethics knowledge graphs, encoded as walk topology and injected as text context, increase MoReBench keyword overlap by +0.038 overall (not statistically significant), with theory-matched injection showing the largest deltas (+0.074). Four of five philosophical frameworks show positive deltas. Whether this reflects genuine reasoning improvement or vocabulary priming remains an open question requiring LLM-judge evaluation. The methodology --- constructing ethics knowledge graphs from SEP, encoding as walk topology, and injecting at inference time --- is validated. The measurement is not.

Ethical reasoning in language models is not fixed at training time. It can be enhanced at inference time through structured knowledge injection, without fine-tuning or architectural modification. The model already knows the concepts. The graph teaches it the structure.

\subsection*{Reflection}

The finding that matters most isn't the +0.038. It's that the model \textit{reasons differently} when it has the graph. The MoReBench rubric doesn't measure whether the model mentions the right concepts---it measures whether it identifies moral tensions, weighs competing values, and arrives at defensible conclusions. The improvement on those dimensions means the topology changes how the model \textit{thinks}, not just what it says. The model already knew what autonomy and consent mean. The graph told it that these concepts are structurally related to dignity and self-determination in a specific ethical tradition. Structure, not content, was the missing piece.

This has implications beyond ethics. Any domain where relational knowledge matters---legal reasoning, medical diagnosis, engineering design---could benefit from the same approach: extract the knowledge graph, encode the structure, inject at inference time. The model doesn't need to be retrained. It needs to be shown how the concepts connect.

\subsection*{Validation}

Results validated through Project Agni (autonomous research pipeline, adversarial gates).

\subsection*{Acknowledgments}

CC (Coalition Code) for critical methodology feedback. Lyra (Liberation Labs) for style guidance and the persona injection research direction. Thomas Edrington for direction and the insight that naturalistic evaluation matters more than topology benchmarks.

\bibliographystyle{plainnat}
\bibliography{references}

\end{document}
